\documentclass[../../main.tex]{subfiles}
\begin{document}

\begin{center}
  {\fontsize{48pt}{58pt}\selectfont \textbf{1962}\normalsize\textbf{年度}}
\end{center}

\begin{table}[htb]
  \centering
  \caption{出題テーマと難易度}
  \begin{tabular}{|c|p{9cm}|c|}
    \hline
    \textbf{問題番号} & \textbf{テーマ} & \textbf{難易度} \\
    \hline
    第1問 & 対数を係数とする二次方程式の根の配置 & \\
    \hline
    第2問 & 直角二等辺三角形内の角度と線分の関係 & \\
    \hline
    第3問 & 傾いた直線上で動く折線の重心 & \\
    \hline
    第4問 & 無限級数の和 & \\
    \hline
    第5問 & 直方体の辺の範囲と体積の最大化 & \\
    \hline
    第6問 & 正弦曲線の接線と囲まれる面積 & \\
    \hline
  \end{tabular}
\end{table}

\begin{table}[htb]
  \centering
  \caption{解答}
  \renewcommand{\arraystretch}{1.4}
  \begin{tabular}{|c|c|p{9cm}|}
    \hline
    \textbf{問題番号} & \textbf{小問} & \textbf{答え} \\
    \hline
    第1問 & -- & $1<a<b$または$b<a<1$ \\
    \hline
    第2問 & -- & $x+y=\sqrt{2}\left(1-\dfrac{\sqrt{3}}{3}\right)$ \\
    \hline
    第3問 & -- & $\tan\theta=3$ \\
    \hline
    第4問 & -- & $|r|>1$のとき$\dfrac{1}{1-r}$，$|r|<1$のとき$\dfrac{r}{1-r}$ \\
    \hline
    第5問 & (1) & $0<x<3$または$3<x\le 4$ \\
    \cline{2-3}
    & (2) & 最大体積は$4$ \\
    \hline
    第6問 & (1) & $R=(\dfrac{14\pi}{3}-2\sqrt{3},\dfrac{4\pi}{3}+2\sqrt{3})$ \\
    \cline{2-3}
    & (2) & $\dfrac{8\pi^2}{3}+8\sqrt{3}\pi-63$ \\
    \hline
  \end{tabular}
\end{table}

\paragraph{試験形態}
6問．試験時間は要確認．

\end{document}
